% ======================================================================
% 第3章 メタヒューリスティクス（最適化手法の基礎理論）
% ======================================================================

% 【追加】章の導入（位置付け）
本章では，本研究で提案する経路制御手法の基礎となる数理最適化理論について述べる．
具体的には，メタヒューリスティクスの概要と，本研究で採用する粒子群最適化 (PSO) の基本アルゴリズム，およびその探索特性について解説する．
これらは第5章で提案する改良手法のベースとなるものである．

\section{概要}
情報通信ネットワークにおける経路制御問題は，制約条件を満たしつつ目的関数を最適化する経路を選択する問題であり，数学的には組合せ最適化問題に帰着される．
特に，本研究で扱う多制約最適経路問題 (MCOP) は，ネットワーク規模の増大に伴い探索空間が指数関数的に広がるため，従来の厳密解法では多項式時間での解の導出が困難である．
このような背景から，自然界の現象や生物の行動を模範としたアルゴリズムであるメタヒューリスティクスが注目されている．
メタヒューリスティクスは，必ずしも厳密な最適解を保証しないものの，現実的な計算時間内で精度の高い近似解を得ることができる手法である．
代表的な手法として，生物の進化を模倣した遺伝的アルゴリズム (GA) や，鳥の群れ行動を模倣した粒子群最適化 (PSO) などが挙げられる．

本章では，本研究で採用する最適化手法であるPSOについて詳述する．
まず最適化問題の基礎概念を定義し，PSOの基本原理とデータ構造，および探索性能を左右するトポロジー構造について解説する．
さらに，PSOが抱える課題である早熟収束と，その対策としてのリスタート戦略について述べる．

\section{最適化問題の基礎}
一般に最適化問題とは，与えられた制約条件の下で，目的関数 $f(x)$ を最小化または最大化する決定変数 $x$ を求める問題と定義される\cite{aiyoshi2002}．
探索空間 $S$ における解 $x \in S$ のうち，全ての $y \in S$ に対して $f(x) \le f(y)$ （最小化問題の場合）を満たす解 $x^*$ を大域的最適解 (Global Optimum) と呼ぶ．
一方，ある近傍領域 $N(x) \subset S$ 内においてのみ最小である解 $x'$ を局所最適解 (Local Optimum) と呼ぶ．
複雑な地形を持つ実問題においては，アルゴリズムが局所最適解に捕らわれ，大域的最適解に到達できない局所解へのスタックが頻繁に発生する．

また，メタヒューリスティクスの設計において最も重要な概念は，探索 (Exploration) と活用 (Exploitation) のバランスである\cite{iba2011}．
探索とは，未知の領域を広く調査し，新たな解の候補を発見する能力であり，大域的探索とも呼ばれる．
一方，活用とは，既に発見した有望な解の周辺を集中的に調査し，精度を高める能力であり，局所探索とも呼ばれる．
探索が不足すれば局所解に陥りやすく，活用が不足すれば収束速度が低下する．
PSOなどの群知能アルゴリズムは，パラメータやトポロジーの調整によってこのバランスを制御する．

\section{粒子群最適化 (PSO)}
PSOは，1995年にKennedyとEberhartによって提案された群知能アルゴリズムである\cite{kennedy1995}．
鳥の群れが餌を探す際の行動ルールをモデル化しており，単純な計算式で高度な探索能力を実現できる特徴がある．
PSOでは，探索空間内を移動する複数の粒子 (Particle) を用いて解を探索する．
各粒子 $i$ は，現在の位置 $x_i$，速度 $v_i$，そして自身がこれまでに発見した最良の位置であるパーソナルベスト $pBest_i$ を保持する．
また，群れ全体または近傍の中で発見された最良の位置をグローバルベスト $gBest$ として共有する．

これらの粒子は，反復 $t$ ごとに以下の式に従って速度と位置を更新しながら探索を行う\cite{Shi1998}．

\begin{equation}
 v_i(t+1) = w \cdot v_i(t) + c_1 r_1 (pBest_i - x_i(t)) + c_2 r_2 (gBest - x_i(t))
 \label{eq:pso_velocity}
\end{equation}

\begin{equation}
 x_i(t+1) = x_i(t) + v_i(t+1)
 \label{eq:pso_position}
\end{equation}

ここで，$w$ は慣性重み (Inertia Weight) であり，前回の速度をどれだけ維持するかという探索の勢いを制御する．
$c_1, c_2$ は加速係数であり，それぞれ自身の経験 ($pBest$) と社会的な情報 ($gBest$) への引き寄せの強さを表す．
$r_1, r_2$ は $[0, 1]$ の範囲の一様乱数であり，探索に確率的な揺らぎを与える．

この更新プロセスの概念を図~\ref{fig:pso_update_concept}に示す．
図中の青色矢印で示される粒子の新しい移動ベクトルは，現在の移動方向を維持しようとする慣性成分，自身が過去に見つけた最良地点へ戻ろうとする認知成分，および群れ全体が見つけた最良地点へ向かおうとする社会成分の3つのベクトルの合成によって決定される．
これらがバランスよく働くことで，粒子は探索空間を効率的に移動し，最適解へと収束していく．

\begin{figure}[tb]
  \centering
  \includegraphics[width=0.7\textwidth]{images/図3-1.pdf}
  \caption{PSOにおける粒子の位置と速度の更新ベクトル図}
  \label{fig:pso_update_concept}
\end{figure}

PSOのアルゴリズム実装において，各粒子は並列的に自身の状態を更新しながら探索を行う．
第 $t$ 世代における $N$ 個の粒子が保持するデータ構造と，情報の流れを図~\ref{fig:pso_particle_table}に示す．
各粒子は位置と速度という物理量に加え，その位置での適応度である評価値を持つ．
各ステップにおいて，粒子は自身の評価値を計算し，過去の自身の最良記録である $pBest$ を更新するか判断する．
さらに，群れ全体の $pBest$ の中から最良のものが $gBest$ として特定され，全粒子に共有される．
この $pBest$ と $gBest$ の情報が，次世代 $t+1$ の速度更新（式\ref{eq:pso_velocity}）に利用されるサイクルを繰り返すことで，解の精度を向上させていく．

\begin{figure}[t]
  \centering
  \includegraphics[width=0.8\textwidth]{images/図3-2.pdf}
  \caption{各世代における粒子のデータ構造と情報の更新}
  \label{fig:pso_particle_table}
\end{figure}

また，粒子間での情報共有の範囲を規定するトポロジー構造は，PSOの探索性能に大きく影響する重要な要素である\cite{Kennedy1999}．
代表的なモデルとして，グローバル型 (gBest Model) とローカル型 (lBest Model) の構造比較を図~\ref{fig:pso_topology}に示す．

図~\ref{fig:pso_topology}左に示すグローバル型は，全ての粒子が互いに情報を共有するスター型トポロジーである．
群れ全体の最良解 $gBest$ に全粒子が急速に引き寄せられるため，収束速度は極めて速いが，局所最適解に陥りやすい（早熟収束）という欠点がある．
一方，図~\ref{fig:pso_topology}右に示すローカル型は，各粒子が隣接する一部の粒子（近傍）とのみ情報を共有するリング型などの構造である．
情報は群れ全体にゆっくりと伝播するため，収束速度は遅くなるが，群れの多様性が維持されやすく，多峰性関数においても局所解を回避して大域最適解を発見する能力が高い\cite{Kennedy1999}．
本研究のような複雑な多制約問題においては，探索の多様性を維持できるlBestモデルの特性が有利に働くと考えられる．

\begin{figure}[b]
  \centering
  \includegraphics[width=0.8\textwidth]{images/図3-3.pdf}
  \caption{gBestトポロジーとlBestトポロジーの構造比較}
  \label{fig:pso_topology}
\end{figure}

\section{早熟収束とリスタート戦略}
PSOの最大の課題は，早熟収束 (Premature Convergence) である\cite{iba2013}．
これは，探索の初期段階で群れ全体が特定の局所最適解に引き寄せられ，粒子間の距離すなわち多様性が失われることで，それ以上の探索が不可能になる現象を指す．

この問題への対策として，本研究ではリスタート戦略 (Restart Strategy) を導入する．
リスタート戦略の概念を図~\ref{fig:restart_concept}に示す．
図~\ref{fig:restart_concept}左は，粒子が局所最適解に集中し，探索が停滞している状態 (Stagnation) を示している．
この状態では，これ以上計算を続けても解の改善は見込めない．

そこで，停滞を検知した場合に，図~\ref{fig:restart_concept}右のように粒子の位置を強制的に再初期化 (Explosion) する．
ただし，これまでに得られた最良の情報（エリート）は保持しつつ，他の粒子を探索空間全体に再配置することで，一度失われた多様性を回復させる．
これにより，アルゴリズムは局所解から脱出し，探索空間の異なる領域から再度大域的最適解を目指すことが可能となる．

\begin{figure}[htb]
  \centering
  \includegraphics[width=1.0\textwidth]{images/図3-4.pdf}
  \caption{早熟収束の発生とリスタートによる再探索の様子}
  \label{fig:restart_concept}
\end{figure}